\section{Conexión con EP139 y EP136}

EP139 aborda la historia técnica de los Agent y destaca el proactive agent y el universal digital agent; EP136 aborda el modelo como OS y destaca que el modelo se convierte en la capa de planificación de tareas; EP135, por su parte, traslada ambos juicios a los productos sociales: si la IA puede completar tareas por iniciativa propia, también puede ayudar por iniciativa propia a que las personas establezcan relaciones. Si el modelo es el OS, la red social podría convertirse en la aplicación con el context personal más sensible sobre ese OS.

\begin{knowledgebox}{Marco de lectura conjunta de los tres episodios}
EP139 ofrece la definición técnica de Agent, EP136 ofrece el panorama industrial del modelo como OS y EP135 ofrece un caso concreto de producto de IA social. En conjunto responden a una pregunta: cómo pasa el modelo de responder preguntas a actuar en nombre de las personas.
\end{knowledgebox}

\subsection{Resumen de la sección}

EP135 extiende los dos episodios anteriores hacia la conversión en producto para el consumidor final. No discute la capacidad del modelo en sí, sino cómo se reescribirán las redes sociales cuando el modelo pueda actuar por iniciativa propia.
